\chapter{Resumen}\label{Chap:Resumen}

El tema de investigación abordado en este trabajo se centra en el diseño de un controlador con doble efecto integral para sistemas inestables con fenómenos de atascamiento-deslizamiento. Los resultados de esta investigación se desarrollaron específicamente en el contexto del sistema de levitación magnética en su configuración inestable. La representación del modelo se basó en los parámetros de la literatura, los cuales fueron obtenidos a través de análisis en experimentos de laboratorio relacionados con este dispositivo. Mediante simulaciones digitales, se llevó a cabo un estudio de la dinámica no lineal asociada a la fricción de Coulomb, la zona muerta (que resulta de la combinación de estos dos fenómenos y se conoce como efecto de atascamiento-deslizamiento) y la fuerza electromagnética, más la inestabilidad del sistema debido a la presencia de múltiples puntos de equilibrio. Se procedió a la linealización del modelo tanto en su configuración estable como inestable, diseñando controladores lineales correspondientes. Estos controladores fueron adaptados a medida que se incorporaban los efectos de deslizamiento-atascamiento y la fricción no lineal en el modelo. Finalmente, se llevaron a cabo pruebas de seguimiento y regulación para validar la efectividad de los controladores diseñados en la estabilización del sistema.